% Lösungen Trigonometrie (zusammenbau v0.4, Heft)
\einheitenkopf{Trigonometrie}

\erg{34}{a) r: G, s: A, t: H \quad b) $\mathrm{tan}\,\gamma = \frac{z}{y}$ \quad c) $\mathrm{sin}\,23^\circ \approx 0{,}391$}
\erg{35}{a) $\mathrm{sin}\,\beta = \frac{e}{f}$ \quad b) $\mathrm{cos}\,\alpha = \frac{n}{m}$ \quad c) $\mathrm{sin}\,\beta = \frac{q}{p}$ \quad d) $\mathrm{cos}\,\alpha = \frac{v}{w}$ \quad e) $\mathrm{tan}\,\gamma = \frac{k}{m}$}
\erg{36}{a) $\mathrm{sin}\,47^\circ = \frac{AB}{840}$, $AB = 840 \cdot \mathrm{sin}\,47^\circ \approx 614$ m \quad b) $AB = 2\,300 \cdot \mathrm{sin}\,41^\circ \approx 1\,509$ m \quad c) $\mathrm{cos}\,53^\circ = \frac{AD}{780}$, $AD = 780 \cdot \mathrm{cos}\,53^\circ \approx 469{,}4$ m \quad d) $AD = 1\,650 \cdot \mathrm{cos}\,62^\circ \approx 774{,}6$ m \quad e) $\mathrm{tan}\,38{,}5^\circ = \frac{1{,}35}{x}$, $x = 1{,}35 : \mathrm{tan}\,38{,}5^\circ \approx 1{,}70$ m \quad f) $x = 0{,}85 : \mathrm{tan}\,27{,}4^\circ \approx 1{,}64$ m \quad g) $\mathrm{sin}\,57^\circ = \frac{6{,}3}{a}$, $a = 6{,}3 : \mathrm{sin}\,57^\circ \approx 7{,}5$ m \quad h) $a = 4{,}8 : \mathrm{sin}\,63^\circ \approx 5{,}4$ cm \quad i) $x = 4 : \mathrm{cos}\,60^\circ = 4 : 0{,}5 = 8$ \quad j) $x = 6 : \mathrm{tan}\,45^\circ = 6 : 1 = 6$ \quad k) $\mathrm{sin}\,\gamma = \frac{d}{f}$ \quad l) $\mathrm{sin}\,\beta = \frac{x}{z}$ \quad m) $\mathrm{tan}\,\alpha = \frac{q}{p}$ \quad n) $\mathrm{tan}\,\gamma = \frac{u}{v}$}
\erg{37}{a) Winkel: Umkehrtaste \quad b) $\mathrm{sin}\,\alpha = \frac{3}{8} = 0{,}375$, $\alpha = \mathrm{sin}^{-1}(0{,}375) \approx 22{,}0^\circ$ \quad c) $\mathrm{cos}\,\alpha = \frac{5}{7}$, $\alpha = \mathrm{cos}^{-1}\left(\frac{5}{7}\right) \approx 44{,}4^\circ$}
\erg{38}{a) $\mathrm{cos}\,\alpha = \frac{460}{610}$, $\alpha \approx 41{,}1^\circ$ \quad b) $\mathrm{cos}\,\alpha = \frac{190}{245}$, $\alpha \approx 39{,}1^\circ$ \quad c) $\mathrm{cos}\,\alpha = \frac{3}{8}$, $\alpha \approx 68{,}0^\circ$, also rund $68^\circ$ \quad d) $\mathrm{cos}\,\alpha = \frac{7}{15}$, $\alpha \approx 62{,}2^\circ$, also rund $62^\circ$ \quad e) $\mathrm{cos}\,\varepsilon = \frac{9}{26} \approx 0{,}346$, $\varepsilon \approx 69{,}7^\circ$, also rund $70^\circ$ \quad f) $\mathrm{cos}\,\varepsilon = \frac{7}{18}$, $\varepsilon \approx 67{,}1^\circ$ \quad g) Teildreieck mit der Höhe: $\mathrm{sin}\,\alpha = \frac{5{,}8}{9{,}6}$, $\alpha \approx 37{,}2^\circ$ \quad h) $\mathrm{sin}\,\alpha = \frac{3{,}9}{8{,}5}$, $\alpha \approx 27{,}3^\circ$ \quad i) $x = \sqrt{240^2 + 18^2} \approx 240{,}7$ cm; $\mathrm{tan}\,\alpha = \frac{18}{240}$, $\alpha \approx 4{,}3^\circ$ \quad j) $x = \sqrt{300^2 + 25^2} \approx 301{,}0$ cm; $\mathrm{tan}\,\alpha = \frac{25}{300}$, $\alpha \approx 4{,}8^\circ$ \quad k) $AB = 4$, $AC = 5$; $BC = \sqrt{4^2 + 5^2} \approx 6{,}40$; $\mathrm{tan}\,\beta = \frac{5}{4}$, $\beta \approx 51{,}3^\circ$ \quad l) $AB = 3$, $AC = 7$; $BC = \sqrt{3^2 + 7^2} \approx 7{,}62$; $\mathrm{tan}\,\beta = \frac{7}{3}$, $\beta \approx 66{,}8^\circ$}
\erg{39}{a) Teildreieck aus Höhe, linkem Schenkel und linkem Überstand: Schenkel H, Höhe G, Überstand A \quad b) Teildreieck: Schenkel H, Höhe G. $\mathrm{sin}\,63^\circ = \frac{h}{7}$, $h = 7 \cdot \mathrm{sin}\,63^\circ \approx 6{,}2$ cm \quad c) $s = 4 : \mathrm{sin}\,62^\circ \approx 4{,}5$ m}
\erg{40}{a) $96{,}5 \cdot \mathrm{sin}\,27^\circ \approx 43{,}81$ m; Hügel $\approx 45{,}2$ m; Oberkante $\approx 54{,}2$ m \quad b) $148 \cdot \mathrm{sin}\,24^\circ \approx 60{,}20$ m; Hügel $\approx 61{,}9$ m; Oberkante $\approx 69{,}4$ m \quad c) Hilfsdreieck: Katheten $18 - 7 = 11$ (an $\alpha$) und $16$; $\mathrm{tan}\,\alpha = \frac{16}{11}$, $\alpha \approx 55{,}5^\circ$ \quad d) Hilfsdreieck: Katheten $13 - 5 = 8$ (an $\alpha$) und $10$; $\mathrm{tan}\,\alpha = \frac{10}{8}$, $\alpha \approx 51{,}3^\circ$ \quad e) $BF = 96 \cdot \mathrm{sin}\,41^\circ \approx 62{,}98$ cm, $AF = 96 \cdot \mathrm{cos}\,41^\circ \approx 72{,}45$ cm; im Dreieck $CBF$: $BC = BF : \mathrm{cos}\,57^\circ \approx 115{,}6$ cm \quad f) $BF = 86 \cdot \mathrm{sin}\,33^\circ \approx 46{,}84$ cm, $AF = 86 \cdot \mathrm{cos}\,33^\circ \approx 72{,}13$ cm; im Dreieck $CBF$: $BC = BF : \mathrm{cos}\,62^\circ \approx 99{,}8$ cm \quad g) Dreieck $AFC$, rechter Winkel bei $F$: $\mathrm{sin}\,19^\circ = \frac{h}{AC}$, $AC = 21{,}4 : \mathrm{sin}\,19^\circ \approx 65{,}7$ cm \quad h) Dreieck $AFC$, rechter Winkel bei $F$: $\mathrm{sin}\,23^\circ = \frac{h}{AC}$, $AC = 15{,}6 : \mathrm{sin}\,23^\circ \approx 39{,}9$ cm \quad i) $\angle CAD = 90^\circ - 59^\circ = 31^\circ$, $\angle DAB = 72^\circ - 31^\circ = 41^\circ$; $BD = 610 \cdot \mathrm{tan}\,41^\circ \approx 530$ m \quad j) $\angle CAD = 90^\circ - 63^\circ = 27^\circ$, $\angle DAB = 71^\circ - 27^\circ = 44^\circ$; $BD = 380 \cdot \mathrm{tan}\,44^\circ \approx 367$ m \quad k) Gebraucht: $AD$, $\delta$ und der rechte Winkel bei $E$. Reihenfolge: $AE = AD \cdot \mathrm{sin}\,\delta$, $DE = AD \cdot \mathrm{cos}\,\delta$, Flächeninhalt = die Hälfte von $AE \cdot DE$. \quad l) Gebraucht: $AD$, $\alpha$ und der rechte Winkel bei $F$. Reihenfolge: $DF = AD \cdot \mathrm{sin}\,\alpha$, $AF = AD \cdot \mathrm{cos}\,\alpha$, Flächeninhalt = die Hälfte von $AF \cdot DF$.}
\erg{41}{a) sin, cos, tan \quad b) $\frac{b}{\mathrm{sin}\,51^\circ} = \frac{7}{\mathrm{sin}\,68^\circ}$, $b = 7 \cdot \mathrm{sin}\,51^\circ : \mathrm{sin}\,68^\circ \approx 5{,}9$ cm \quad c) $PQ$ liegt gegenüber $R$, $PR$ gegenüber $Q$: $PR = 8 \cdot \mathrm{sin}\,53^\circ : \mathrm{sin}\,69^\circ \approx 6{,}8$ cm}
\erg{42}{a) $BD$ liegt gegenüber $118^\circ$, $AD$ gegenüber $33^\circ$: $BD = 4{,}6 \cdot \mathrm{sin}\,118^\circ : \mathrm{sin}\,33^\circ \approx 7{,}5$ cm \quad b) $GS = 3\,200 \cdot \mathrm{sin}\,57^\circ : \mathrm{sin}\,72^\circ \approx 2\,822$ m; Weg $\approx 4\,622$ m $\approx 4{,}6$ km}
\erg{43}{a) $\gamma = 180^\circ - 31^\circ - 113^\circ = 36^\circ$; $BC = 520 \cdot \mathrm{sin}\,31^\circ : \mathrm{sin}\,36^\circ \approx 455{,}6$ m \quad b) $\gamma = 180^\circ - 27^\circ - 121^\circ = 32^\circ$; $BC = 290 \cdot \mathrm{sin}\,27^\circ : \mathrm{sin}\,32^\circ \approx 248{,}4$ m \quad c) dritter Winkel $180^\circ - 6^\circ - 137^\circ = 37^\circ$ (gegenüber $220$ cm); $y = 220 \cdot \mathrm{sin}\,137^\circ : \mathrm{sin}\,37^\circ \approx 249{,}3$ cm \quad d) dritter Winkel $180^\circ - 4^\circ - 148^\circ = 28^\circ$ (gegenüber $190$ cm); $y = 190 \cdot \mathrm{sin}\,148^\circ : \mathrm{sin}\,28^\circ \approx 214{,}5$ cm \quad e) Das Dreieck ist nicht rechtwinklig: Sinussatz. $AC = 9 \cdot \mathrm{sin}\,62^\circ : \mathrm{sin}\,4^\circ \approx 113{,}92$ m, also rund $114$ m \quad f) Das Dreieck ist nicht rechtwinklig: Sinussatz. $AC = 8 \cdot \mathrm{sin}\,51^\circ : \mathrm{sin}\,6^\circ \approx 59{,}48$ m, also rund $59{,}5$ m \quad g) $\angle ACB = 31^\circ$, $\angle ACD = 57^\circ$; $AD = 7 \cdot \mathrm{sin}\,57^\circ : \mathrm{sin}\,78^\circ \approx 6{,}0$ m \quad h) $\angle ACB = 37^\circ$, $\angle ACD = 59^\circ$; $AD = 10 \cdot \mathrm{sin}\,59^\circ : \mathrm{sin}\,69^\circ \approx 9{,}2$ m \quad i) Winkel bei $D$: $39^\circ$; $DE = 9 \cdot \mathrm{sin}\,118^\circ : \mathrm{sin}\,39^\circ \approx 12{,}63$ m; $DP = DE - 8 \approx 4{,}6$ m \quad j) Winkel bei $D$: $46^\circ$; $DE = 16 \cdot \mathrm{sin}\,107^\circ : \mathrm{sin}\,46^\circ \approx 21{,}27$ m; $DP = DE - 17 \approx 4{,}3$ m \quad k) $\angle ABD = 42^\circ$, $\angle DBC = 62^\circ$; $BC = 1\,300 \cdot \mathrm{sin}\,57^\circ : \mathrm{sin}\,62^\circ \approx 1\,235$ m – nicht gleich $1\,300$ m, die Aussage ist falsch. \quad l) $\angle ABD = 39^\circ$, $\angle DBC = 73^\circ$; $BC = 1\,760 \cdot \mathrm{sin}\,63^\circ : \mathrm{sin}\,73^\circ \approx 1\,640$ m – nicht gleich $1\,760$ m, die Aussage ist falsch.}
